\documentclass[11pt,letterpaper]{article}
\usepackage[utf8]{inputenc}
\usepackage{amsmath}
\usepackage{amsfonts}
\usepackage{amssymb}
\usepackage{geometry}
\usepackage{hyperref}
\usepackage{xcolor}

\definecolor{mpdark}{HTML}{0F172A}
\definecolor{mpblue}{HTML}{3B82F6}
\definecolor{mpgreen}{HTML}{10B981}

\geometry{margin=1in}

\title{\vspace{-20pt}\textbf{\color{mpdark}Money Plot Labs Technical Whitepaper 002 \\ The Asymmetric Glide Path: Solving Bifurcated Growth}}
\author{\textbf{Money Plot Labs}}
\date{June 2026}

\begin{document}

\maketitle

\section{The Transcendental Challenge}
In Technical Whitepaper 001, we derived a closed-form solution for the working years $w$ under a unified growth rate $r$. However, modern portfolio theory suggests a "glide path," where risk—and therefore expected return—is reduced as an individual moves from accumulation to decumulation. 

When we introduce two distinct rates, $r_{acc}$ (Working Growth) and $r_{dec}$ (Retirement Growth), the boundary equation becomes:
\begin{equation}
    R_f(A_w(w, r_{acc}), f, r_{dec}) = L
\end{equation}

Expanding this using the Annuity Due model derived previously:
\begin{equation}
    \left[ A_0(1+r_{acc})^w + s \frac{(1+r_{acc})^w - 1}{r_{acc}} \right] (1+r_{dec})^{H-w} - c \frac{(1+r_{dec})^{H-w+1} - (1+r_{dec})}{r_{dec}} = L
\end{equation}

Unlike the unified case, the variable $w$ now appears in the exponents of two different bases: $(1+r_{acc})$ and $(1+r_{dec})$. This is a \textbf{transcendental equation}. There is no finite sequence of algebraic operations that can isolate $w$, necessitating a shift from analytical solvers to numerical approximation.

\section{The Numerical Solver: Bisection Method}
To solve for $w$, the Money Plot engine utilizes the Bisection Method (a form of Binary Search). This method is robust, guaranteed to converge on a continuous function, and avoids the "division by zero" or "log of negative" errors that Newton-Raphson methods might encounter in extreme financial scenarios.

\subsection{The Objective Function}
We define the function $f(w)$ as the net difference between the terminal balance and the target legacy floor:
\begin{equation}
    f(w) = \text{TerminalBalance}(w) - L
\end{equation}
The goal is to find the root where $f(w) = 0$ within the interval $[0, H]$.

\subsection{The Algorithm}
\begin{enumerate}
    \item \textbf{Initialize}: Set $low = 0$, $high = H$.
    \item \textbf{Evaluate Boundaries}: 
        \begin{itemize}
            \item If $f(0) \geq 0$, the individual is "Already Free" ($w=0$).
            \item If $f(H) < 0$, the target is unreachable ("Never").
        \end{itemize}
    \item \textbf{Iterate}: For $i = 1$ to 60:
        \begin{itemize}
            \item $mid = \frac{low + high}{2}$
            \item If $f(mid) < 0$, then $low = mid$ (Must work longer).
            \item Else, $high = mid$ (Can retire sooner).
        \end{itemize}
    \item \textbf{Result}: $w = \frac{low + high}{2}$.
\end{enumerate}
A 60-iteration search provides a precision of $H \cdot 2^{-60}$. Over a 100-year horizon, this yields a precision of approximately $8.6 \times 10^{-17}$ years, far exceeding any practical requirement.

\section{Bifurcated Model Components}

\subsection{Accumulation (The Input)}
The asset state at the moment of retirement ($A_w$) is governed by the accumulation rate:
\begin{equation}
    A_w = A_0(1+r_{acc})^w + s \frac{(1+r_{acc})^w - 1}{r_{acc}}
\end{equation}

\subsection{Decumulation (The Output)}
The depletion of $A_w$ over the retirement years $f = H-w$ is governed by the decumulation rate using the Annuity Due timing (spending at the start of the year):
\begin{equation}
    Balance_t = A_w(1+r_{dec})^t - c \frac{(1+r_{dec})^{t+1} - (1+r_{dec})}{r_{dec}}
\end{equation}

\section{Impact of Rate Bifurcation}
Splitting the rates reveals two critical insights:
\begin{enumerate}
    \item \textbf{Accumulation Dominance}: Higher $r_{acc}$ reduces $w$ exponentially. A 1\% increase in $r_{acc}$ during early working years has a larger impact on the retirement date than a 1\% increase in $r_{dec}$ during retirement.
    \item \textbf{Sequence Protection}: Lowering $r_{dec}$ (e.g., shifting to bonds) increases the required $A_w$ but flattens the depletion curve, reducing the variance of the terminal outcome $L$.
\end{enumerate}

\end{document}
